% ch24_revisions_texsis.tex
%
% Suggested Revisions to Chapter 24 of Recursive Macroeconomic Theory
% Plain TeX using the TeXsis macro package (version 2.18 or later).
%
% To compile:   texsis ch24_revisions_texsis
% Output:       ch24_revisions_texsis.dvi
%
% TeXsis is available from https://www.texsis.org or CTAN:
%   macros/plain/formats/texsis
%
% ---------------------------------------------------------------

\texsis          % initialize TeXsis: 12pt, doublespaced, paper format
\input paper           % set up for a research paper / preprint

\title{Suggested Revisions to Chapter~24 of {\sl Recursive Macroeconomic
Theory\/}}

\author{Based on Barth\'elemy and Mengus (2020)}

\abstract
This note collects suggested revisions to the statements and proofs in
Chapter~24 of Ljungqvist and Sargent (2018), following the correction
note of Barth\'elemy and Mengus~(2020).  The core issue is that the
two-value characterization of supported equilibria stated on page~1034
is presented as a necessary and sufficient condition but is in fact only
sufficient at the point where it is invoked.  It becomes necessary only
{\it after\/} the set of equilibrium payoffs~$V$ is known to be compact,
which is precisely what the surrounding argument aims to establish.  The
fix requires replacing each admissible {\it 4-tuple\/} $(x,y,w_1,w_2)$
with an admissible {\it triple\/} $(x,y,w(\cdot))$ where $w(\cdot)$ is
a function on~$Y$, proving compactness of~$V$ with this refined
machinery, and then recovering the two-value formulation as a corollary.
The economic content of the chapter is entirely unaffected.
\endabstract

\body           % output title block and begin body text

% ---------------------------------------------------------------
\section{The Gap and Its Source}
% ---------------------------------------------------------------

On page~1034, the text states that a competitive equilibrium
$(x,y)\in C$ is supported by a subgame perfect equilibrium (SPE)---that
is, $x=\sigma^h_1$ and $y=\sigma^g_1$---{\it if and only if\/} there
exist two values $(v_1,v_2)\in V\times V$ such that
%
$$  (1-\delta)\,r(x,y)+\delta v_1
    \;\geq\;
    (1-\delta)\,r(x,\eta)+\delta v_2 ,
    \qquad \forall\,\eta\in Y .
    \eqno{(1)}  $$
%
No explicit proof of the ``only if'' direction is provided.  As we show
below, condition~(1) is {\it sufficient\/} but not {\it necessary\/}
unless $V$ is already known to be compact---but compactness of~$V$ is
precisely what the subsequent argument aims to establish.

The correct necessary and sufficient condition replaces the two scalar
continuation values $(v_1,v_2)$ with a {\it function\/}
$v(\cdot):Y\to V$, one continuation value for each possible government
action.  We state and prove this refined lemma in
Section~2, use it to rebuild the APS machinery in Sections~3 and~4,
prove compactness of~$V$ in Section~5, and then recover the two-value
characterization as a corollary in Section~6.

% ---------------------------------------------------------------
\section{Revised Characterization of Supported Equilibria}
% ---------------------------------------------------------------

The following lemma replaces the ``if and only if'' assertion on
page~1034.

\Lemma  {\bf (Necessary and sufficient condition for support by an
SPE.)}
A competitive equilibrium $(x,y)\in C$ is supported by an SPE, i.e.\
$x=\sigma^h_1$ and $y=\sigma^g_1$, if and only if there exists a
function $v(\cdot):Y\to V$ such that
%
$$  (1-\delta)\,r(x,y)+\delta v(y)
    \;\geq\;
    (1-\delta)\,r(x,\eta)+\delta v(\eta) ,
    \qquad \forall\,\eta\in Y .
    \eqno{(2)}  $$
\endLemma

\Proof
{\sl Necessity.\/}
Suppose $(x,y)\in C$ is supported by an SPE~$\sigma$, so that
$x=\sigma^h_1$ and $y=\sigma^g_1$.
By condition~(1) of the Lemma on page~1030, for each $\eta\in Y$ the
continuation strategy $\sigma|_{(x,\eta)}$ is itself an SPE.
Define $v(\eta)=V_g(\sigma|_{(x,\eta)})\in V$ for each $\eta\in Y$.
Condition~(3) of the Lemma on page~1030 then gives
$$  (1-\delta)\,r(x,y)+\delta V_g(\sigma|_{(x,y)})
    \;\geq\;
    (1-\delta)\,r(x,\eta)+\delta V_g(\sigma|_{(x,\eta)}) ,  $$
which is exactly~(2) with $v(\eta)=V_g(\sigma|_{(x,\eta)})$.

\smallskip
{\sl Sufficiency.\/}
Suppose there exists $v(\cdot):Y\to V$ satisfying~(2).
Since $v(\eta)\in V$ for each $\eta\in Y$, there exists for each~$\eta$
an SPE $\tilde\sigma(\eta)$ with $V_g(\tilde\sigma(\eta))=v(\eta)$.
Construct a strategy profile~$\sigma$ by setting
$\sigma|_{(x,\eta)}=\tilde\sigma(\eta)$ for all $\eta\in Y$, and
(without loss of generality) $\sigma|_{(\nu,\eta)}=\tilde\sigma(\eta)$
for all $\nu\in X$.  By construction this profile satisfies all three
conditions of the Lemma on page~1030 and is therefore an SPE with
$\sigma^h_1=x$ and $\sigma^g_1=y$.
\endProof

\medskip
Condition~(1) of the original text is {\it sufficient\/}: given
$(v_1,v_2)\in V\times V$, set $v(y)=v_1$ and $v(\eta)=v_2$ for
$\eta\neq y$, and~(2) holds.  The reverse implication holds once~$V$
is known to be compact (see Section~6), but cannot be assumed at this
stage.

% ---------------------------------------------------------------
\section{Revised Definition~7: Admissible Triples}
% ---------------------------------------------------------------

Using the Lemma of Section~2, Definition~7 on page~1034 should be
replaced as follows.

\Definition  {\bf (Admissible triple.)}
Let $W\subset{\bf R}$.  A triple $(x,y,w(\cdot))$ is said to be
{\it admissible with respect to\/}~$W$ if $(x,y)\in C$ and
$w(\cdot):Y\to W$ is a function satisfying
$$  (1-\delta)\,r(x,y)+\delta w(y)
    \;\geq\;
    (1-\delta)\,r(x,\eta)+\delta w(\eta) ,
    \qquad \forall\,\eta\in Y .  $$
\endDefinition

% ---------------------------------------------------------------
\section{Revised Definition~8 and the Monotonicity of~$B$}
% ---------------------------------------------------------------

Definition~8 on page~1036 is updated to reflect the triple-based notion
of admissibility.

\Definition  {\bf (The operator~$B$.)}
For each set $W\subset{\bf R}$, let $B(W)$ be the set of possible
values
$$  w \;=\; (1-\delta)\,r(x,y)+\delta\,w(y)  $$
associated with admissible triples $(x,y,w(\cdot))$ with respect
to~$W$.
\endDefinition

The monotonicity property is unchanged in statement; only its proof is
adjusted.

\medskip\noindent
{\bf Property (Monotonicity of~$B$).}
{\sl If $W\subseteq W'\subseteq{\bf R}$, then $B(W)\subseteq B(W')$.}

\Proof
Suppose $w\in B(W)$.  Then there exists an admissible triple
$(x,y,w(\cdot))$ with respect to~$W$ such that
$w=(1-\delta)r(x,y)+\delta w(y)$.  Since $W\subseteq W'$, the function
$w(\cdot)$ maps~$Y$ into~$W'$ as well, so $(x,y,w(\cdot))$ is
admissible with respect to~$W'$ and $w\in B(W')$.
\endProof

% ---------------------------------------------------------------
\section{Revised Theorem~1 and Compactness of~$V$}
% ---------------------------------------------------------------

Definition~9 (self-generating sets) is unchanged: $W$ is
self-generating if $W\subseteq B(W)$.  Theorem~1 is restated and its
proof adapted to use admissible triples.

\Theorem  {\bf (Self-generating sets are subsets of~$V$.)}
If $W\subset{\bf R}$ is bounded and self-generating, then
$B(W)\subseteq V$.
\endTheorem

\Proof
Assume $W\subseteq B(W)$.  Choose $w\in B(W)$ and construct an SPE
with value~$w$ as follows.

\medskip\noindent
{\it Step~1.\/}
Because $w\in B(W)$, there exist $(x_1,y_1)\in C$ and a function
$w_1(\cdot):Y\to W$ such that
$$  w \;=\; (1-\delta)\,r(x_1,y_1)+\delta\,w_1(y_1)
    \;\geq\;
    (1-\delta)\,r(x_1,\eta)+\delta\,w_1(\eta) ,
    \quad \forall\,\eta\in Y .  $$
Set $\sigma_1=(x_1,y_1)$.

\medskip\noindent
{\it Step~2.\/}
Since $w_1(y_1)\in W\subseteq B(W)$, there exist $(x_2,y_2)\in C$ and
a function $w_2(\cdot):Y\to W$ such that
$$  w_1(y_1) \;=\; (1-\delta)\,r(x_2,y_2)+\delta\,w_2(y_2)
    \;\geq\;
    (1-\delta)\,r(x_2,\eta)+\delta\,w_2(\eta) ,
    \quad \forall\,\eta\in Y .  $$
Set the first-period outcome of the continuation economy (occurring
given that $y_1$ was chosen in period~1) equal to $(x_2,y_2)$; that
is, set $(\sigma|_{(x_1,y_1)})_1=(x_2,y_2)$.

\medskip\noindent
{\it Continuation.\/}
Iterating in this way, for each $w\in B(W)$ we construct sequences of
functions $w_1(\cdot),w_2(\cdot),w_3(\cdot),\ldots$ and
first-period outcomes $(x_1,y_1),(x_2,y_2),(x_3,y_3),\ldots$.
Since $\delta<1$ and $W$ is bounded,
$$  w
    \;=\; \lim_{t\to\infty}
           \left[
             (1-\delta)\sum_{s=1}^{t}\delta^{s-1}r(x_s,y_s)
             +\delta^t w_t(y_t)
           \right]
    \;=\; (1-\delta)\sum_{s=1}^{\infty}\delta^{s-1}r(x_s,y_s) .  $$
At each stage the continuation values make one-period deviations
unimprovable, so the constructed strategy profile is an SPE.  By
construction $V_g(\sigma)=w$, so $w\in V$.
\endProof

GGHH

\medskip
With the revised machinery in place, compactness of~$V$ follows by an
argument parallel to the proof of Theorem~4 in Abreu, Pearce, and
Stacchetti~(1990).

\Proposition  {\bf (Compactness of~$V$.)}
The set~$V$ is compact.
\endProposition

\Proof
We know that $V$ is bounded and self-generating ($V\subseteq B(V)$).
Let ${\rm cl}(V)$ denote the closure of~$V$.  Since $V$ is bounded,
${\rm cl}(V)$ is compact and $V\subseteq{\rm cl}(V)$.

By monotonicity of~$B$,
$$  V \;=\; B(V) \;\subseteq\; B({\rm cl}(V)) .  $$
Since $B(\cdot)$ maps compact sets into compact sets,
$B({\rm cl}(V))$ is compact.  Because ${\rm cl}(V)$ is the smallest
closed set containing~$V$, and $V\subseteq B({\rm cl}(V))$, it follows
that
$$  {\rm cl}(V) \;\subseteq\; B({\rm cl}(V)) ,  $$
so ${\rm cl}(V)$ is bounded and self-generating.  Theorem~1 then gives
${\rm cl}(V)\subseteq V$, and therefore $V={\rm cl}(V)$, which means
$V$ is compact.
\endProof

% ---------------------------------------------------------------
\section{Corollary: The Two-Value Characterization}
% ---------------------------------------------------------------

Now that $V$ is known to be compact, the function-valued condition of
the Lemma in Section~2 collapses to the economically transparent
two-value form, and the assertion on page~1034 can be stated cleanly
as the following corollary.

\Corollary  {\bf (Two-value characterization.)}
A competitive equilibrium $(x,y)\in C$ is supported by an SPE, that
is $x=\sigma^h_1$ and $y=\sigma^g_1$, if and only if there exist
$(v_1,v_2)\in V\times V$ such that
$$  (1-\delta)\,r(x,y)+\delta v_1
    \;\geq\;
    (1-\delta)\,r(x,\eta)+\delta v_2 ,
    \qquad \forall\,\eta\in Y .  $$
\endCorollary

\Proof
{\sl Sufficiency\/} is immediate: given $(v_1,v_2)\in V\times V$,
define $v(y)=v_1$ and $v(\eta)=v_2$ for $\eta\neq y$; then~(2) holds,
so the Lemma in Section~2 applies.

\medskip
{\sl Necessity.\/}
Since $(x,y)$ is supported by an SPE, the Lemma in Section~2 yields a
function $v(\cdot):Y\to V$ satisfying~(2).  Set $v_1=v(y)$.  The
set $\{v(\eta):\eta\in Y\}$ is a subset of~$V$; since $V$ is compact
(Proposition of Section~5) and bounded, its infimum is attained in~$V$.
Set
$$  v_2 \;=\; \inf_{\eta\in Y}\,v(\eta) \;\in\; V .  $$
Then $v(\eta)\geq v_2$ for all $\eta\in Y$, so~(2) implies
$$  (1-\delta)\,r(x,y)+\delta v_1
    \;=\; (1-\delta)\,r(x,y)+\delta v(y)
    \;\geq\; (1-\delta)\,r(x,\eta)+\delta v(\eta)
    \;\geq\; (1-\delta)\,r(x,\eta)+\delta v_2 ,  $$
for all $\eta\in Y$, as required.
\endProof

% ---------------------------------------------------------------
\section{Summary of Changes}
% ---------------------------------------------------------------

The following list gives the page-level replacements needed in
Chapter~24.

\medskip
\item{$\bullet$}
Page~1034, main ``if and only if'' assertion: replace with the Lemma
of Section~2 above.

\item{$\bullet$}
Page~1034, Definition~7 (admissible 4-tuple): replace with the
Definition of Section~3 above (admissible triple).

\item{$\bullet$}
Page~1036, Definition~8 ($B(W)$ via 4-tuples): replace with the
Definition of Section~4 above.

\item{$\bullet$}
Page~1036, Monotonicity proof (uses $(w_1,w_2)$ pair): replace with
the proof in Section~4 above.

\item{$\bullet$}
Page~1037, Theorem~1 proof (uses scalar $w_1,w_2$): replace with the
proof in Section~5 above (uses functions $w_t(\cdot)$).

\item{$\bullet$}
Page~1038, after Theorem~1 (compactness asserted without proof): add
the Proposition and proof of Section~5 above.

\item{$\bullet$}
Page~1034, two-value assertion: demote to a Corollary, stated and
proved as in Section~6 above (using compactness of~$V$).

\medskip\noindent
All other material in Chapter~24---the economic motivation, the
examples (infinite repetition of Nash, trigger strategies, when
reversion to Nash is not bad enough), the iteration algorithm for
computing~$V$, and the discussion of best and worst SPE values---is
unaffected by these changes.

% ---------------------------------------------------------------
\section*{References}
% ---------------------------------------------------------------

\medskip\noindent\hangindent=1.5em
Abreu, D., D.~Pearce, and E.~Stacchetti (1990): ``Toward a Theory of
Discounted Repeated Games with Imperfect Monitoring,''
{\sl Econometrica\/}, 58, 1041--1063.

\medskip\noindent\hangindent=1.5em
Barth\'elemy, J.\ and E.~Mengus (2020): ``Note on Chapter~24 on
Credible Government Policies: Proof of the Compactness of the Set of
Equilibrium Payoffs,'' unpublished manuscript, Banque de France and
HEC Paris.

\medskip\noindent\hangindent=1.5em
Ljungqvist, L.\ and T.~J.~Sargent (2018): {\sl Recursive
Macroeconomic Theory}, 4th~ed., MIT Press.

\bye
